\section{Introduction}
\label{sec:intro}




Large Language Model (LLM) agents have achieved remarkable success in real-world interactive tasks, including writing complex code~\citep{claudecode,code2math}, operating computers~\citep{UI-TARS,UI-TARS-2}, and supporting scientific discovery~\citep{DR-Tulu,Deep-Research-Agents}.
While these applications differ in domain, they share a common structure: agent capability depends on sustained interaction throughout task execution~\citep{ReAct,park2023generativeagentsinteractivesimulacra}.
This interaction typically unfolds along two closely linked dimensions: \textbf{interaction with users}, through which agents infer goals and preferences, and \textbf{interaction with the external world}, where they gather information and take actions via tools and interfaces~\citep{xi2023risepotentiallargelanguage,Wang_2024}.
Because such interaction is inherently multi-step, it usually requires \textbf{planning}~\citep{huang2024understandingplanningllmagents}: agents must anticipate future outcomes~\citep{current-agent-fail-wm-as-tool} and dynamically adapt their actions as interaction unfolds~\citep{costbench} across both their engagement with users and the external world.
Planning therefore naturally spans both dimensions, reflecting the dual structure of agent interaction. 
% \xiusi{Add figure 1 as a motivating figure to distinguish the roles of user and external world.}

However, real-world planning is rarely unconstrained. Because interaction has a dual structure, agents must handle \textbf{dual constraints} from both the user and the world: user constraints such as preferences and priorities, and world constraints such as tool availability and resource limitations.
Existing benchmarks typically consider only one side of this problem, focusing on either user constraints~\citep{UserBench,opentom,URS-dataset} or world constraints~\citep{appworld,tau-2-bench,PlanBench}, leaving their joint handling largely unexplored.
This raises our central research question: \textbf{can LLM agents plan effectively under both user and world constraints?} 
In practice, this problem is further complicated by two key challenges:
(1) \textbf{Progressive constraint disclosure}: constraints are often implicit rather than specified upfront, requiring agents to uncover them incrementally through proactive exploration.
(2) \textbf{Large action and solution spaces}: real-world tasks involve vast spaces of possible actions and solutions, making performance harder to measure. % ~\citep{newell1972human}
These challenges call for a rigorous benchmark that evaluates whether agents can adaptively plan under dual constraints that are progressively revealed in open-ended planning settings.

\begin{table*}[!t]
% \vspace{-10mm}
\centering
\resizebox{\linewidth}{!}{
    \begin{tabular}{lcccccccc}
    \toprule
    \textbf{Benchmark} 
    & \textbf{\makecell{Iterative\\Re-planning}} 
    & \textbf{\makecell{User\\Interaction}}
    & \textbf{\makecell{World\\Interaction}}
    & \textbf{\makecell{Dual\\Constraint}} 
    & \textbf{\makecell{Progressive\\Disclosure}}
    & \textbf{\makecell{Open-Ended\\ Evaluation}} 
    & \textbf{\makecell{Scalable\\Constraints}} \\
    \midrule
    % Benchmarks rows go here
    \textit{CostBench}~\citep{costbench} & \gmark & \xmark & \cmark & \gmark & \xmark & \xmark & \gmark
    \\
    \textit{FlowBench}~\citep{FlowBench} & \gmark & \cmark & \cmark & \gmark & \gmark & \xmark & \xmark
    \\ % \gmark for Progressive Disclosure: user preference are progressively disclosed, but env constraints are not (workflows are upfront)
    % \xmark for open evaluation: rely on workflow and ground-truth sessions
    % \gmark for dual constraints: The workflow itself is predefined, and users only make demands within this framework. So if you mean "user constraint" you mean: "User can define a set of constraints by himself, and agent must satisfy two independent sets of constraints, workflow and user." FlowBench does not support this design.
    \textit{NaturalPlan}~\citep{NATURAL-PLAN} & \xmark & \xmark & \xmark & \cmark & \xmark & \gmark & \cmark
    \\ % \xmark for Dynamic State, User Interaction, World Interaction, Progressive Disclosure: Static QA task, not agentic task.
    \textit{PrefEval}~\citep{PrefEval} & \xmark & \cmark & \xmark & \xmark & \xmark & \gmark & \cmark
    \\ % \xmark for Progressive Disclosure: "Preference revelation" occurs solely during the dataset construction phase, not during the actual benchmark execution. When the benchmark is run, the model's input consists of a pre-constructed, complete dialogue history that already incorporates the pre-designed preference revelation process. From the model's perspective, the context is effectively "frozen"; no new preferences suddenly emerge during runtime.
    \textit{RealPref}~\citep{RealPref} & \xmark & \xmark & \xmark & \xmark & \cmark & \cmark & \xmark
    \\ % Only user prefs. \xmark for the initial four: During the training and building phases, the dataset consists of multi-turn conversations (with multiple turns per session). However, during the evaluation phase, the agent does not engage in actual multi-turn dialogue; instead, it simply receives a conversation history and responds to a standalone test query.
    \textit{UserBench}~\citep{UserBench} & \xmark & \cmark & \cmark & \xmark & \cmark & \cmark & \cmark
    \\ % \xmark for Dynamic State: It possesses a certain degree of dynamism, yet it is not event-driven. UserBench supports "search failure simulation"—meaning that search results can be corrupted by noise according to a configured probability—to simulate real-world network or tool errors. However, this constitutes pre-configured, controlled noise, rather than unexpected events generated randomly at runtime.
    % \xmark for Dual Constraint: Only cost-effectiveness for world constraints, that doesn't count.
    \textit{PersonaMem-v2}~\citep{PersonaMem-v2} & \xmark & \cmark & \xmark & \xmark & \cmark & \cmark & \cmark
    \\ % No world constraint
    $\tau$-Bench~\citep{tau-bench} & \xmark & \cmark & \cmark & \gmark & \gmark & \gmark & \xmark
    \\ 
    % \gmark on Progressive Disclosure: the policy is known from the beginning to the agent, only the user prefs needs to be inferred from runtime.
    % \gmark on open evaluation: action space and solution space is not inifinite.
    % \gmark for dual constraints: domain policy are provided upfront
    \textit{$\tau^2$-Bench}~\citep{tau-2-bench} & \xmark & \cmark & \cmark & \xmark & \cmark & \xmark & \cmark
    \\ % \xmark for open eval: has strict process and outcome check for eval
    % \xmark on dual constraint: no user preference
    % \cmark for Progressive Disclosure: environment state is progressively disclosed through user's mouth
    \textit{TravelPlanner}~\citep{TravelPlanner} & \xmark & \xmark & \cmark & \cmark & \xmark & \cmark & \cmark
    \\ % \xmark for open eval: has strict process and outcome check for eval
    % \xmark on dual constraint: no user preference
    % \cmark for Progressive Disclosure: environment state is progressively disclosed through user's mouthwi
    \midrule
    \textit{\BENCH{} \textbf{(Ours)}} & \cmark & \cmark & \cmark & \cmark & \cmark & \cmark & \cmark
    \\
    \bottomrule
    \end{tabular}
}
\vspace{-1.5mm}
\caption{Comparison of \bench{} with prior related benchmarks across seven key properties. For each benchmark, the table reports whether each trait is fully (\cmark), partially (\gmark), or not (\xmark) addressed. 
Detailed explanations are provided in Appendix~\ref{app:traits-exp}.}
\label{tab:comparison-table}
\vspace{-0.2in}
\end{table*}

To this end, we introduce \BENCH{}, a dynamic and interactive benchmark evaluating LLM's adaptive planning ability under world and user constraints. 
% \bench{} is rooted in the household domain, where both user preferences and environmental constraints arise naturally in everyday scenarios~\citep{VirtualHome}.
We build \bench{} on top of the MacGyver dataset~\citep{MacGyver}, starting from a curated subset of 307 household-domain instances.
Using a scalable automated pipeline, we augment each task with world constraints that capture environmental limitations and user constraints that capture grounded personal preferences.
This setup allows any solution that satisfies these constraints, while preserving a large and effectively unbounded action space.
During evaluation, constraints are withheld at the outset and disclosed progressively when the agent proposes violating actions, forcing the agent to iteratively and adaptively re-plan in response to newly revealed constraints.

We evaluate ten leading open-source and proprietary LLMs on \bench{} and find that even the strongest model achieves only 67.75\% accuracy, while open-weight models typically remain at or below 30\%.
Moreover, planning quality becomes harder to sustain as progressively disclosed constraints accumulate along a trajectory, and deteriorates further as the overall constraint burden increases.
These challenges are not easily alleviated by explicit constraint tracking or rubric-based feedback, and are particularly severe when user constraints account for a large share of the difficulty.
Further analysis suggests that, under accumulated dual constraints, planning failures are often marked by reduced goal effectiveness and weaker physical plausibility.
Overall, \bench{} lays a foundation for future research on adaptive planning agents, motivating the development of models that can dynamically plan, adapt, and revise under real-world dual constraints.
% \vspace{-0.2in}
% Additionally, when the agent cannot directly access the environment, the user becomes the only source from which the agent can obtain feedback~\citep{}. 
% In such cases, users often exhibit a crucial characteristic: progressive information disclosure~\citep{UserBench}.


